\chapter*{Abstract}
\addcontentsline{toc}{chapter}{Abstract}

Speech enhancement in real-world environments remains challenging due to diverse noise types, non-stationary interference, and limited computational resources on edge devices. This report presents \projectshort{}, an end-to-end intelligent audio denoising analysis platform that integrates data-driven scene characterization, adaptive mathematical denoising, lightweight deep-learning refinement, and interactive evaluation.

The proposed pipeline consists of three layers: (1) a full-dimensional scene analysis module that extracts time--frequency statistics and derives parameter guidance for the Adaptive Noise Adaptive Spectral Subtraction (ANASS) framework; (2) a library of seven composable denoising algorithms centered on an enhanced OM-LSA/MCRA core, orchestrated by a scene-aware rule router; and (3) a dilated-residual TinyResidualRefiner student network (\texttt{student\_runK.pt}, 56\,129 parameters) trained via knowledge distillation from DeepFilterNet3 to perform residual mask correction on mathematically denoised outputs.

Experiments on a 15-scene benchmark covering white, pink, hum, hiss, chirp, reverberation, and babble-like noise demonstrate consistent improvements over fixed-parameter baselines. When optional runK distillation refinement is enabled, all 15 evaluated scenes show reduced RMS error relative to routed mathematical denoising alone (mean improvement $1.39\times10^{-3}$; mean SNR 12.69\,dB $\to$ 14.06\,dB). The web platform exposes runK via \textbf{auto} routing plus \textbf{Run distill refine}.

\vspace{1em}
\noindent\textbf{Keywords:} speech enhancement, spectral subtraction, OM-LSA, MCRA, adaptive routing, knowledge distillation, audio denoising, scene analysis
